\centerline{\bf \Huge 6 класс}
\problem{6.1} На доске написаны две суммы:\par
1 + 22 + 333 + 4444 + 55555 + 666666 +7777777 + 88888888 + 999999999 \par
9 + 98 + 987 + 9876 + 98765 + 987654 + 9876543 + 98765432 + 987654321 \\ 
Определите, какая из них больше (или они равны).

\problem{6.2} На какую цифру оканчивается число $2023^{2022}$?

\problem{6.3} В Древней Спарте выживали сильнейшие. Поэтому 300 спартанцев решили выяснить, кто из них достойный воин. Они разбились по парам и сделали тренировочный бой. Потом опять разбились по парам (возможно, уже по-другому) и снова провели тренировочный бой. Тот, кто выигрывал оба боя, считался достойным воином, а остальных сбрасывали со скалы. Какое минимальное количество воинов могло остаться в Спарте? \par
Известно, что все 300 спартанцев были разной силы, а выигрывал в бою всегда тот, у кого силы больше.

%https://problems.ru/view_problem_details_new.php?id=64440

\problem{6.4} Лера, Аркадий и Владимир записали на доске по одному числу $-$ 0, 2 и 5 соответственно. Спустя минуту каждый из них одновременно стирает свое число и записывает вместо него сумму двух других(то есть Лера запишет 7, Аркадий 5, а Владимир 2). Спустя ещё минуту они повторяют это действие и так несколько раз. Могла ли сумма чисел Леры, Аркадия и Владимира в какой-то момент стать равной 20232022?\\

\problem{6.5} Очир выписал на доску числа $1, \dfrac{1}{2}, \dfrac{1}{3}, \dfrac{1}{4}, \dotsc, \dfrac{1}{2023}$. За один ход он может стереть любые 
\\
\\
два числа $a$ и $b$ и заменить их на число $ab+a+b$. Какое число останется, если Очир проделает 2022 такие операции?

%https://problems.ru/view_problem_details_new.php?id=35225